% GENERATED FILE. Edit canonical lessons, then run npm run generate:books.
% canonical-source-hash: fnv1a64:fa0aa4d0eb175447
% canonical-lessons: GE-C36-mein, GE-C36-dein, GE-C36-sein-his, GE-C36-ihr-her, GE-C36-das-ist, GE-R36-wer-ist-das

\chapter{Mine, Yours, His, Hers}
\label{ch:ge-possessiv}
\hlchaptermodality{49}

\begin{chapteropening}
\textbf{By the end of this chapter:} \emph{I can say whose something is with all four possessives, and introduce somebody who is not me.}

Every person and every object the book owns is claimed and then done something to. The four possessives cost their first two letters and nothing else, because they wear the endings ein bought, kein banked and den confirmed -- six words, one ending set, and the reader had already paid for it three times.
\end{chapteropening}

\section[mein]{mein — the word you have been saying since you first gave your name}

\label{lesson:GE-C36-mein}



\begin{quote}
Say \textbf{mein Name ist …}. You have said it since the second chapter as a single block. The first word of it is about to become a word of its own.

\begin{itemize}
\raggedright
  \item \emph{Recall:} two chapters back, the phrase that says a thing exists — \emph{es gibt}
\end{itemize}
\end{quote}

\subsection*{You'll want to know: mein}
\begin{quote}
\textbf{mein} — \textquotedblleft{}my\textquotedblright{}.
\end{quote}

\begin{quote}
\textbf{mein Bruder} · \textbf{meine Schwester} · \textbf{mein Brot}
\end{quote}

English splits this word in two — \emph{my} in front of a noun, \emph{mine} standing alone — and the split is late and unnecessary. \textbf{mein} and \textbf{mine} are one word, Proto-Germanic \emph{\textbf{*mīnaz}}, and German never divided it.

\begin{grammarlens}[title={Grammar Lens: the family of six}]
\begin{quote}
\textbf{mein} takes the endings \textbf{ein} takes. So does \textbf{kein}, and so will the three words in the next three lessons.
\end{quote}

You paid for this two chapters ago without being told. \emph{Ein} becomes \textbf{eine} before a \emph{die}-word and \textbf{einen} when a \emph{der}-word is done to. \emph{Kein} does the same. \textbf{Mein} does the same:

\begin{quote}
\textbf{mein Bruder} $\to$ \textbf{Ich habe meinen Bruder.} \textbf{meine Schwester} $\to$ \textbf{Ich habe meine Schwester.}
\end{quote}

One ending set, six words wearing it. Learning the fourth, fifth and sixth costs nothing at all, because the endings were bought with \emph{ein}, banked again with \emph{kein}, and confirmed a third time with \emph{den}.
\end{grammarlens}

\subsection*{Your turn}
Say the article first, so the ending falls out of it.

Say these aloud:

\begin{itemize}
\raggedright
  \item \emph{der} Bruder $\to$ \textquotedblleft{}mein Bruder\textquotedblright{}, then done-to, \textquotedblleft{}meinen Bruder\textquotedblright{}
  \item \emph{die} Schwester $\to$ \textquotedblleft{}meine Schwester\textquotedblright{}, which does not move
  \item \textquotedblleft{}meine Mutter, mein Vater\textquotedblright{}
  \item \textquotedblleft{}mein Freund, meine Freundin\textquotedblright{}
  \item the three that move on the same row — \textquotedblleft{}einen, keinen, meinen\textquotedblright{}
\end{itemize}

\emph{Twice through:} \textquotedblleft{}mein Bruder — meinen Bruder.\textquotedblright{}

\begin{tcolorbox}[breakable,colback=teal!4,colframe=teal!35!black,title={Before you move on}]
Which endings does \emph{mein} take? (\textbf{The ones \emph{ein} takes} — and \emph{kein}.) Say \textquotedblleft{}my brother\textquotedblright{} as the thing you have. (\textbf{Meinen Bruder}.) Which two English words are \emph{mein}? (\emph{\textbf{My}} and \emph{\textbf{mine}}, split late and needlessly.)
\end{tcolorbox}
\section[dein]{dein — \textquotedblleft{}your,\textquotedblright{} and the third \emph{th} that became a \emph{d}}

\label{lesson:GE-C36-dein}



\begin{quote}
Say \textbf{du} and \textbf{dir}. Both were once English \emph{thou} and \emph{thee}. Predict the German for \emph{thine} before you read on.
\end{quote}

\subsection*{You'll want to know: dein}
\begin{quote}
\textbf{dein} — \textquotedblleft{}your\textquotedblright{}, to somebody you call \emph{du}.
\end{quote}

\begin{quote}
\textbf{dein Bruder} · \textbf{deine Schwester} · \textbf{deinen Bruder}
\end{quote}

The endings are \emph{mein}'s endings, which were \emph{ein}'s endings. Nothing new is being learned here except the first two letters.

\begin{cousinweb}
Three words, one shift, and you have now seen all three:

\noindent
\begin{tabularx}{\linewidth}{@{}>{\raggedright\arraybackslash}X>{\raggedright\arraybackslash}X@{}}
\toprule
\textbf{English} & \textbf{German} \\
\midrule
\textbf{th}ou & \textbf{d}u \\
\textbf{th}ee & \textbf{d}ir \\
\textbf{th}ine & \textbf{d}ein \\
\bottomrule
\end{tabularx}

Germanic \emph{\textbf{*þ}} hardened into German \textbf{d} across the board, so any English word beginning \emph{th-} is worth testing against a German \emph{d-}. \emph{That} is \textbf{das}. \emph{Three} is \textbf{drei}. \emph{Thank} is \textbf{danken}. The rule has been paying out since your first article and it will keep paying.
\end{cousinweb}

\subsection*{Your turn}
Say these aloud:

\begin{itemize}
\raggedright
  \item \textquotedblleft{}dein Bruder, deine Schwester\textquotedblright{}
  \item the pair — \textquotedblleft{}mein Freund, dein Freund\textquotedblright{}
  \item a question with it — \textquotedblleft{}Wer ist deine Freundin?\textquotedblright{}
  \item the \emph{th} to \emph{d} rule, four times — \textquotedblleft{}that/das, three/drei, thank/danke, thine/dein\textquotedblright{}
\end{itemize}

\emph{Twice through:} \textquotedblleft{}Wer ist dein Bruder?\textquotedblright{}

\begin{tcolorbox}[breakable,colback=teal!4,colframe=teal!35!black,title={Before you move on}]
Say \textquotedblleft{}your sister\textquotedblright{} and \textquotedblleft{}your brother\textquotedblright{}. Which English word is \emph{dein}? (\emph{\textbf{Thine}}.) Which sound law lines the two up? (\textbf{Germanic \emph{þ} becoming German \emph{d}.}) Whose endings does \emph{dein} wear? (\emph{\textbf{Ein\emph{'s\textbf{ — like \emph{kein} and \emph{mein}.)}}}}
\end{tcolorbox}
\section[sein]{sein — \textquotedblleft{}his,\textquotedblright{} and a collision worth meeting head-on}

\label{lesson:GE-C36-sein-his}



\begin{quote}
\textbf{Sein} is the verb \emph{to be}, the one with three roots braided into it. Say \textbf{er ist}. Now meet a second, unrelated word spelled the same way.
\end{quote}

\subsection*{You'll want to know: sein}
\begin{quote}
\textbf{sein} — \textquotedblleft{}his\textquotedblright{}.
\end{quote}

\begin{quote}
\textbf{sein Bruder} · \textbf{seine Schwester} · \textbf{seinen Bruder}
\end{quote}

Same endings again — \emph{ein}'s. Third word in a row that costs nothing but its first letters.

The two \emph{sein}s are genuinely two words. The possessive is Proto-Germanic \emph{\textbf{*sīnaz}}, a cousin of Latin \textbf{suus}. The verb is the braided one. They arrived at the same spelling from opposite directions and German never separated them.

\begin{grammarlens}[title={Grammar Lens: how a reader tells them apart}]
The position settles it, every time, without exception:

\begin{quote}
\textbf{sein Bruder} — a noun follows, so this is \emph{his}. \textbf{Ich will sein} — no noun follows, so this is \emph{to be}.
\end{quote}

A possessive is always leaning on a noun. An infinitive is not. German has more of these collisions than English does — \emph{sie} alone runs three jobs — and each time, what resolves it is the shape of the sentence around the word rather than the word itself.
\end{grammarlens}

\subsection*{Your turn}
Say these aloud:

\begin{itemize}
\raggedright
  \item \textquotedblleft{}sein Bruder, seine Schwester\textquotedblright{}
  \item the three so far — \textquotedblleft{}mein Vater, dein Vater, sein Vater\textquotedblright{}
  \item \textquotedblleft{}Das ist sein Hund.\textquotedblright{}
  \item the collision, both ways — \textquotedblleft{}sein Bruder\textquotedblright{} and the bare verb, \textquotedblleft{}sein\textquotedblright{}
\end{itemize}

\emph{Twice through:} \textquotedblleft{}mein Vater, dein Vater, sein Vater.\textquotedblright{}

\begin{tcolorbox}[breakable,colback=teal!4,colframe=teal!35!black,title={Before you move on}]
Say \textquotedblleft{}his mother\textquotedblright{} and \textquotedblleft{}his brother\textquotedblright{}. What else is the word \emph{sein}? (\textbf{The verb \emph{to be}.}) What tells them apart? (\textbf{Whether a noun follows.}) Say the three possessives you now own. (\emph{\textbf{Mein, dein, sein}}.)
\end{tcolorbox}
\section[ihr]{ihr — \textquotedblleft{}her,\textquotedblright{} and a word with three jobs}

\label{lesson:GE-C36-ihr-her}



\begin{quote}
\textbf{Ihr} already means \emph{you all}. \textbf{Sie} already means \emph{she}, \emph{they} and formal \emph{you}. German is comfortable with a word doing several jobs, and this lesson adds one more.
\end{quote}

\subsection*{You'll want to know: ihr}
\begin{quote}
\textbf{ihr} — \textquotedblleft{}her\textquotedblright{}.
\end{quote}

\begin{quote}
\textbf{ihr Bruder} · \textbf{ihre Schwester} · \textbf{ihren Bruder}
\end{quote}

Said \emph{eer}: the \textbf{h} is silent and stretches the \emph{i}, exactly as it does in \emph{Ihnen}. The endings are \emph{ein}'s, for the fourth time.

\begin{grammarlens}[title={Grammar Lens: three jobs, all settled by position}]
Where the word stands says which job it is doing:

\begin{itemize}
\raggedright
  \item \textbf{ihr wohnt} — a verb with \emph{-t} follows, so this is \emph{you all}.
  \item \textbf{ihr Bruder} — a noun follows, so this is \emph{her}.
  \item \textbf{Ihr Bruder}, capitalised — the formal \emph{your}, which the capital marks the same way it marks formal \emph{Sie}.
\end{itemize}

Nothing here is ambiguous in a real sentence, and the capital does real work — a third reason for the capitalisation rule the writing chapter taught.

The same \emph{ihr} also means \emph{their}, which is one job too many to add today; the book will come back to it.
\end{grammarlens}

\subsection*{Your turn}
Say these aloud:

\begin{itemize}
\raggedright
  \item \textquotedblleft{}ihr Bruder, ihre Schwester\textquotedblright{}
  \item all four — \textquotedblleft{}mein, dein, sein, ihr\textquotedblright{}
  \item the two jobs side by side — \textquotedblleft{}ihr wohnt\textquotedblright{} and \textquotedblleft{}ihr Hund\textquotedblright{}
  \item \textquotedblleft{}Das ist ihre Mutter.\textquotedblright{}
\end{itemize}

\emph{Twice through:} \textquotedblleft{}mein, dein, sein, ihr.\textquotedblright{}

\begin{tcolorbox}[breakable,colback=teal!4,colframe=teal!35!black,title={Before you move on}]
Name the four possessives. (\emph{\textbf{Mein, dein, sein, ihr}}.) What tells you which job \emph{ihr} is doing? (\textbf{What follows it} — a verb or a noun.) What does the capital on \emph{Ihr} mean? (\textbf{The formal \emph{your}}.)
\end{tcolorbox}
\section[Das ist mein Bruder]{das ist — introducing somebody who is not you}

\label{lesson:GE-C36-das-ist}



\begin{quote}
You can give your own name three ways and shake hands afterwards. You cannot yet introduce anybody else, and the phrase that does it is two words you already have.
\end{quote}

\subsection*{You'll want to know: das ist}
\begin{quote}
\textbf{das ist} — \textquotedblleft{}this is\textquotedblright{}, \textquotedblleft{}that is\textquotedblright{}.
\end{quote}

\begin{quote}
\textbf{Das ist mein Bruder.} — \textquotedblleft{}This is my brother.\textquotedblright{}
\end{quote}

\begin{quote}
\textbf{Das ist meine Schwester.} — \textquotedblleft{}This is my sister.\textquotedblright{}
\end{quote}

And the reply is one you have owned since your first introduction: \textbf{Freut mich.}

\begin{grammarlens}[title={Grammar Lens: the das that ignores gender}]
Look at the second sentence again. \textbf{Schwester} is a \emph{die}-word, and the phrase still opens with \textbf{das}. That is not a mistake, and it is not the article.

The \textbf{das} at the front of the sentence is the older, pointing word the article descends from — the demonstrative behind English \textbf{that}. It points at a person without committing to a gender, so it stays \textbf{das} in front of everybody:

\begin{quote}
\textbf{Das ist mein Vater.} · \textbf{Das ist meine Mutter.} · \textbf{Das ist mein Hund.}
\end{quote}

English does the same thing and hides it as well: \emph{that is my sister} uses a word that started life as the neuter \emph{þæt} and long ago stopped caring.
\end{grammarlens}

\subsection*{Your turn}
Say these aloud:

\begin{itemize}
\raggedright
  \item \textquotedblleft{}Das ist mein Bruder.\textquotedblright{} — then answer, \textquotedblleft{}Freut mich.\textquotedblright{}
  \item \textquotedblleft{}Das ist meine Schwester.\textquotedblright{} — the \emph{das} does not become \emph{die}
  \item the household — \textquotedblleft{}Das ist meine Mutter. Das ist mein Vater.\textquotedblright{}
  \item with the other three possessives — \textquotedblleft{}Das ist dein Freund. Das ist seine Freundin. Das ist ihr Hund.\textquotedblright{}
\end{itemize}

\emph{Twice through:} \textquotedblleft{}Das ist meine Schwester. — Freut mich.\textquotedblright{}

\begin{tcolorbox}[breakable,colback=teal!4,colframe=teal!35!black,title={Before you move on}]
Introduce your sister. (\textbf{Das ist meine Schwester.}) Why is it not \emph{die ist}? (\textbf{That \emph{das} is the old pointing word, not the article}, and it ignores gender.) What do you say back? (\textbf{Freut mich.})
\end{tcolorbox}
\section[Practice]{Wer ist das? — the chapter, run cold}

\label{lesson:GE-R36-wer-ist-das}



\begin{quote}
Name the four possessives. Then name the two older words whose endings all four borrow, and say why that means the four cost almost nothing.
\end{quote}

\subsection*{Your turn: introduce everybody}
Every person this book has named, introduced and claimed.

Say these aloud:

\begin{itemize}
\raggedright
  \item \textquotedblleft{}Das ist mein Bruder. Das ist meine Schwester.\textquotedblright{}
  \item \textquotedblleft{}Das ist meine Mutter. Das ist mein Vater.\textquotedblright{}
  \item the two collective words — \textquotedblleft{}meine Eltern, meine Geschwister\textquotedblright{}
  \item the pair the \emph{-in} suffix built — \textquotedblleft{}mein Freund, meine Freundin\textquotedblright{}
  \item \textquotedblleft{}Das ist meine Familie.\textquotedblright{} then answer, \textquotedblleft{}Freut mich.\textquotedblright{}
  \item the four Latin cousins those words keep — \emph{frater, soror, mater, pater}
\end{itemize}

\subsection*{Your turn: the distant band — everything, claimed}
Put \emph{mein} in front of everything, and then do something to it.

Say these aloud:

\begin{itemize}
\raggedright
  \item the animals — \textquotedblleft{}mein Hund, meinen Hund; meine Katze\textquotedblright{}
  \item the table — \textquotedblleft{}mein Kaffee, meinen Kaffee; mein Brot; meine Milch\textquotedblright{}
  \item the body, all of it — \textquotedblleft{}mein Kopf, meine Hand, mein Arm, mein Fuß, mein Herz, mein Auge, mein Ohr, mein Mund, meine Nase\textquotedblright{}
  \item the two abstractions — \textquotedblleft{}mein Name\textquotedblright{}, \textquotedblleft{}meine Zeit\textquotedblright{}
  \item the six words on one ending set — \textquotedblleft{}ein, kein, mein, dein, sein, ihr\textquotedblright{}
  \item the \emph{der}-row, five words deep — \textquotedblleft{}den, einen, keinen, meinen, deinen\textquotedblright{}
\end{itemize}

\subsection*{Your turn: the far band — the last three chapters, over one family}
Everything since the negation chapter, said about one family.

Say these aloud:

\begin{itemize}
\raggedright
  \item joined and corrected — \textquotedblleft{}Das ist nicht mein Bruder, sondern mein Vater.\textquotedblright{} \textquotedblleft{}Das ist meine Schwester, aber sie ist müde.\textquotedblright{}
  \item with a reason both ways — \textquotedblleft{}Das ist mein Bruder, denn er hat meinen Namen.\textquotedblright{} \textquotedblleft{}Ich frage meine Mutter, weil ich keine Zeit habe.\textquotedblright{}
  \item existence and its marked form — \textquotedblleft{}Es gibt keinen Kaffee.\textquotedblright{} \textquotedblleft{}Es gibt einen Bruder oder eine Schwester.\textquotedblright{}
  \item the three W-words, aimed at people — \textquotedblleft{}Wer ist das? Wann? Warum?\textquotedblright{}
  \item the two words built out of \emph{not} and \emph{one} — German's \textbf{kein} and English's \textbf{none}
  \item the two W-words whose English twins begin \emph{wh-} — \textquotedblleft{}wer/who, wann/when\textquotedblright{}
\end{itemize}

\begin{tcolorbox}[breakable,colback=teal!4,colframe=teal!35!black,title={Before you move on}]
Say all six words that share one ending set. (\emph{\textbf{Ein, kein, mein, dein, sein, ihr}}.) Which row of endings moves? (\textbf{The \emph{der} row.}) Introduce your sister and answer yourself. (\textbf{Das ist meine Schwester. — Freut mich.}) Which other job does \emph{sein} do, and which other two does \emph{ihr} do?
\end{tcolorbox}
